\chapter{A Working Landscape of Conformal Field Theories}
\label{ch:volume-iii-landscape-cfts}

\section{Purpose of the Survey}

The preceding chapters described a conformal field theory by local data:
primary operators, their conformal representations, OPE coefficients, stress
tensor Ward identities, and crossing consistency.  The present chapter records
standard families of examples and constructions.  Its role is organizational:
it locates the abstract data in concrete classes of theories before the
two-dimensional Virasoro analysis begins.

Throughout this chapter, a local CFT means a unitary Poincare-invariant
Lorentzian theory, or its Euclidean continuation, with a stress tensor and a
state-operator correspondence on the sphere.  Other conformal systems, such as
one-dimensional conformal quantum mechanics, use related representation theory
but have a different locality structure.

\section{Free Conformal Theories}

The free massless scalar in \(D\) dimensions has a conformal stress tensor
after the improvement discussed earlier.  Its fundamental scalar has
\[
  \Delta_\phi=\frac{D-2}{2},
\]
and saturates the scalar unitarity bound.  The equation
\[
  \partial^2\phi=0
\]
is the null-state shortening condition in its conformal multiplet.

Free fermions also give conformal theories in any dimension in which the
spinor representation is defined.  A free Dirac fermion has dimension
\[
  \Delta_\psi=\frac{D-1}{2},
\]
and the Dirac equation is the corresponding shortening condition.

The free Maxwell field is conformal in \(D=4\).  Its stress tensor is traceless
in four dimensions, and the gauge-invariant field strength \(F_{\mu\nu}\) has
dimension \(2\).  In dimensions other than four, the Maxwell action remains
scale covariant with a dimensionful normalization, while the local conformal
stress-tensor structure differs from the four-dimensional one.

\section{Two Dimensions}

In two dimensions, local conformal transformations are holomorphic and
antiholomorphic maps.  The stress tensor decomposes into two chiral currents,
and their modes generate the Virasoro algebra.  This infinite-dimensional
symmetry produces families with much sharper algebraic control than is
available in a generic higher-dimensional CFT.

Important two-dimensional classes include:
\begin{itemize}
  \item Virasoro minimal models, including the Ising and tricritical Ising
  fixed points;
  \item Wess--Zumino--Witten models, whose chiral algebras are affine Lie
  algebras;
  \item coset models, obtained by quotienting chiral current algebras;
  \item compact bosons and toroidal theories, including orbifolds;
  \item Liouville and Toda theories, which have continuous spectra;
  \item nonlinear sigma models whose conformality is encoded by beta-function
  equations for target-space data.
\end{itemize}
The rest of this volume develops the Virasoro and modular structures needed
to make these statements precise.

\section{Three Dimensions}

Three-dimensional CFTs arise naturally as critical points of statistical
systems and as infrared fixed points of quantum field theories.  The critical
Ising CFT is the fixed point of a \(\mathbb Z_2\)-invariant scalar theory.
The \(O(N)\) vector models are fixed points with global \(O(N)\) symmetry,
whose large-\(N\) expansion gives a controlled approximation to their operator
data.

Fermionic fixed points include Gross--Neveu and Gross--Neveu--Yukawa theories.
Gauge theories provide further examples: Chern--Simons--matter theories are
conformal when the Chern--Simons level is quantized and matter interactions
are tuned to a fixed point; QED\(_3\) with sufficiently many light fermions or
scalars has candidate conformal regimes.  Supersymmetry produces many
additional examples with protected sectors and exact constraints.

The numerical conformal bootstrap is especially effective in three
dimensions.  For the Ising CFT, the bootstrap kink and mixed-correlator
islands isolate a small region in the space of CFT data, including the leading
\(\mathbb Z_2\)-odd scalar \(\sigma\), the leading even scalar
\(\epsilon\), and the stress-tensor coefficient \(C_T\).

\section{Four Dimensions}

Four dimensions contain weakly coupled fixed points, strongly coupled fixed
points, and supersymmetric families with exactly marginal deformations.

The Banks--Zaks construction starts from \(SU(N_c)\) gauge theory with
\(N_f\) massless fundamental Dirac fermions.  When \(N_f\) lies just below the
loss of asymptotic freedom, the two-loop beta function has a perturbative
zero.  In the large-\(N_c\) estimate, the conformal window begins near
\begin{equation}
  \frac{34N_c^3}{13N_c^2-3}
  \lesssim
  N_f
  \lesssim
  \frac{11}{2}N_c .
  \label{eq:banks-zaks-window-estimate}
\end{equation}
The upper endpoint is where asymptotic freedom is lost; the lower estimate is
where the two-loop fixed point ceases to be weakly coupled.

Supersymmetric four-dimensional CFTs include \(\mathcal N=1\) gauge theories
with nontrivial infrared fixed points, \(\mathcal N=2\) theories with Coulomb
and Higgs branch structures, and \(\mathcal N=4\) super Yang--Mills theory.
The last has an exactly marginal complex coupling and will be discussed later
as the most symmetric four-dimensional gauge-theory example.

\section{Five and Six Dimensions}

Five-dimensional superconformal field theories can have relevant
deformations to supersymmetric gauge theories.  In this situation the gauge
theory is an effective description on part of the moduli space, while the UV
fixed point is specified by local CFT data and protected observables.

Six-dimensional CFTs include \((2,0)\) theories labeled by simply laced Lie
algebras and a broad class of \((1,0)\) theories.  They are often constructed
through string-theoretic or geometric methods, and their compactifications
generate many lower-dimensional theories.  Their existence illustrates that
the local CFT framework is broader than a Lagrangian presentation.

\section{Higher Dimensions and Classification}

Complete classification of unitary local CFTs above two dimensions remains an open problem.
The conformal bootstrap gives rigorous exclusion methods and numerical
islands under specified assumptions, while perturbation theory, large-\(N\)
expansions, supersymmetry, and geometric constructions give complementary
families of examples.  The monograph uses all these constructions as sources
of data and checks, while keeping the logical definition of a CFT tied to
local operators, states, Ward identities, OPE, and crossing.
